\documentclass{article}

\usepackage[T1]{fontenc}
\usepackage[utf8]{inputenc}
\usepackage{microtype}
\usepackage{graphicx}
\usepackage{subcaption}
\usepackage{booktabs}
\usepackage{multirow}
\usepackage{hyperref}

\newcommand{\theHalgorithm}{\arabic{algorithm}}

\usepackage{icml2026}

\usepackage{amsmath}
\usepackage{amssymb}
\usepackage{mathtools}
\usepackage{amsthm}

\usepackage[capitalize,noabbrev]{cleveref}

% Defaults; codex_numbers.tex (generated by build_short_codex.py)
% \renewcommand's these once the runs finish.
\providecommand{\BestSchedROWP}{TBD}
\providecommand{\DeltaSchedROWP}{TBD}
\providecommand{\SchedStructuralCharacterization}{TBD}
\providecommand{\FullModalChoice}{TBD}
\providecommand{\NarrowBroadOutcome}{TBD}
\providecommand{\StrengthOfEvidence}{TBD}
\providecommand{\CodexSchedTrainSI}{TBD}
\providecommand{\CodexSchedTrainSII}{TBD}
\providecommand{\CodexSchedTrainSIII}{TBD}
\providecommand{\CodexSchedRowpSI}{TBD}
\providecommand{\CodexSchedRowpSII}{TBD}
\providecommand{\CodexSchedRowpSIII}{TBD}
\providecommand{\CodexSchedStructSI}{TBD}
\providecommand{\CodexSchedStructSII}{TBD}
\providecommand{\CodexSchedStructSIII}{TBD}
\providecommand{\CodexFullTrainSI}{TBD}
\providecommand{\CodexFullTrainSII}{TBD}
\providecommand{\CodexFullTrainSIII}{TBD}
\providecommand{\CodexFullRowpSI}{TBD}
\providecommand{\CodexFullRowpSII}{TBD}
\providecommand{\CodexFullRowpSIII}{TBD}
\providecommand{\CodexFullModalSI}{TBD}
\providecommand{\CodexFullModalSII}{TBD}
\providecommand{\CodexFullModalSIII}{TBD}
\providecommand{\CodexSchedTrainMean}{TBD}
\providecommand{\CodexSchedTrainStd}{TBD}
\providecommand{\CodexSchedRowpMean}{TBD}
\providecommand{\CodexSchedRowpStd}{TBD}
\providecommand{\CodexNarrowVerdict}{TBD}
\providecommand{\CodexBroadVerdict}{TBD}
\IfFileExists{codex_numbers.tex}{\input{codex_numbers}}{}

\icmltitlerunning{FunWake-Codex: Narrow vs broad with OpenAI Codex CLI}

\begin{document}

\twocolumn[
\icmltitle{FunWake-Codex: A Frontier-Agent Replication of\\
the Narrow-Interface Discovery Result}

\icmlsetsymbol{equal}{*}

\begin{icmlauthorlist}
\icmlauthor{Anonymous Authors}{anon}
\end{icmlauthorlist}

\icmlaffiliation{anon}{Anonymous Institution}

\icmlcorrespondingauthor{Anonymous Authors}{anon@example.com}

\icmlkeywords{LLM agents, algorithm discovery, wind farm optimization, benchmark, OpenAI Codex}

\vskip 0.3in
]

\printAffiliationsAndNotice{}

\begin{abstract}
Prior work introduced FunWake, a benchmark that exposes LLM agents to a
narrow (4-output schedule) vs.\ broad (write an arbitrary optimizer)
interface on the same wind-farm-layout problem, and reported that
Anthropic Claude Code and Google Gemini CLI both \emph{discovered novel
schedule parameterizations} given the narrow interface but
\emph{re-implemented community-standard solvers} given the broad one.
We replicate that experiment with a third frontier agentic CLI,
OpenAI Codex (\texttt{gpt-5.5}), running 3 independent agent seeds
in each interface (6 runs total). [PLACEHOLDER]: Codex's
schedule-only runs reach a held-out AEP of \BestSchedROWP\,GWh
($+\DeltaSchedROWP$ vs.\ baseline), structurally
\SchedStructuralCharacterization\ relative to Claude/Gemini's
discoveries.  Codex's full-optimizer runs converge on
\FullModalChoice\ across seeds.  The narrow-vs-broad asymmetry
\NarrowBroadOutcome\ for Codex, providing
\StrengthOfEvidence\ evidence that the FunWake narrow interface
elicits algorithmic novelty across the major frontier agentic CLIs of
2026.
\end{abstract}

\section{Background}

We use FunWake as defined in the companion paper. Briefly: the
benchmark places $N$ wind turbines inside a polygon to maximize annual
energy production (AEP) subject to boundary and minimum-spacing
constraints, using a differentiable Gaussian wake
model~\cite{bastankhah2014new}. Two interfaces are exposed to the
agent:

\paragraph{Narrow.}
The agent writes \texttt{schedule\_fn(step, total\_steps, lr0,
alpha0)} returning \texttt{(lr, alpha, beta1, beta2)}, controlling a
fixed 8000-step Adam skeleton. The skeleton is identical across
agents and submissions, so evaluation cost is independent of the
schedule's complexity.

\paragraph{Broad.}
The agent writes an arbitrary
\texttt{optimize(sim, n\_target, boundary, min\_spacing, wd, ws,
weights)} function. The agent controls everything: initialization,
solver, constraint handling, and stopping.

The training cell is DEI ($N{=}50$, observed wind rose); the held-out
cell is ROWP ($N{=}74$, IEA 10\,MW turbine, Weibull rose). The
500-multi-start TopFarm SGD~\cite{pedersen2019topfarm} baseline
scores 5540.7\,GWh on DEI and 4246.7\,GWh on ROWP.

In the companion paper, Claude Code and Gemini CLI given the narrow
interface produced novel schedule structures (dual Gaussian LR bumps
for Claude, 8-cycle cosine restarts with cyclic Adam betas for
Gemini) that beat the baseline on the held-out cell. Given the broad
interface, both reimplemented SLSQP~\cite{kraft1988software} or
wrapped \texttt{topfarm\_sgd\_solve}.

\section{Codex Method}

We add OpenAI Codex CLI v0.125 (\texttt{gpt-5.5}) as a third
frontier agentic CLI, alongside Claude Code and Gemini CLI. Codex
runs locally via the \texttt{codex exec} non-interactive command
under a workspace-write sandbox. Each agent run uses the same
hot-start (\texttt{results/seed\_schedule.py} for narrow,
\texttt{results/seed\_optimizer.py} for broad), the same
\texttt{agent\_memory.md} cross-iteration scaffolding, the same
attempt-log + scoring tools, and an \texttt{AGENTS.md} task spec
mirroring the per-mode CLAUDE.md/GEMINI.md prompts used for
Claude/Gemini.

\paragraph{Time budget.}
Schedule-only saturation in Claude/Gemini occurred near attempt 150
in the original 5\,hr runs. Using the empirical formula
$t_\mathrm{total} = (1 + \mathrm{think}_\mathrm{factor}) \cdot N
\cdot t_\mathrm{seed}$ with $N{=}150$, $t_\mathrm{seed}\!\approx\!30$\,s,
and $\mathrm{think}_\mathrm{factor}\!\approx\!1$ (empirical), we
allocate 3\,hr per schedule-only run and 4.5\,hr per full-optimizer
run.

\paragraph{Seeds.}
We launch 3 independent agent runs per interface (6 runs total).
Codex CLI single-account concurrency caps the chain to one run at a
time; total wall-clock $\approx 22.5$\,hr.

\paragraph{Scoring.}
For each run we extract the highest-training-AEP feasible script and
score it on the held-out ROWP cell.  We additionally classify each
script's structural motif via regex matching against a fixed taxonomy
(Gaussian bumps, cosine restarts, cyclic betas, alpha-squeeze for
narrow; SLSQP-wrapper, SGD-wrapper, custom-Adam, GA, CMA-ES for
broad).

\section{Results}

\subsection{Schedule-only ($n=3$)}

\begin{table}[t]
\caption{Codex schedule-only across 3 agent seeds. Best feasible
script per run, scored on the held-out ROWP cell. Comparators are
Claude/Gemini single-best from the companion paper.}
\label{tab:codex_sched}
\vskip 0.1in
\begin{center}
\begin{small}
\begin{sc}
\begin{tabular}{lrrr}
\toprule
Run                     & Train  & ROWP    & Structures \\
\midrule
Codex seed 1            & \CodexSchedTrainSI   & \CodexSchedRowpSI    & \CodexSchedStructSI   \\
Codex seed 2            & \CodexSchedTrainSII  & \CodexSchedRowpSII   & \CodexSchedStructSII  \\
Codex seed 3            & \CodexSchedTrainSIII & \CodexSchedRowpSIII  & \CodexSchedStructSIII \\
\midrule
Codex mean              & \CodexSchedTrainMean & \CodexSchedRowpMean & --- \\
\quad $\pm$std          & \CodexSchedTrainStd  & \CodexSchedRowpStd  & ---  \\
\midrule
Claude (companion best) & 5555.7 & 4271.5  & dual-bumps  \\
Gemini (companion best) & 5565.1 & 4269.3  & cos-restart \\
500-start baseline      & 5540.7 & 4246.7  & ---         \\
\bottomrule
\end{tabular}
\end{sc}
\end{small}
\end{center}
\vskip -0.1in
\end{table}

[PLACEHOLDER: 1--2 paragraphs describing the Codex schedule-only
results from \cref{tab:codex_sched}. Key questions to address:
(a) does Codex beat the baseline on the held-out cell?
(b) is the variance across 3 seeds comparable to Claude/Gemini's
init-variance?
(c) does Codex rediscover Claude's bump motif, Gemini's cosine-restart
motif, or a third independent structure?]

\subsection{Full-optimizer ($n=3$)}

\begin{table}[t]
\caption{Codex full-optimizer across 3 agent seeds. ``Modality''
classifies the high-level algorithmic approach.}
\label{tab:codex_full}
\vskip 0.1in
\begin{center}
\begin{small}
\begin{sc}
\begin{tabular}{lrrl}
\toprule
Run                     & Train  & ROWP    & Modality \\
\midrule
Codex seed 1            & \CodexFullTrainSI   & \CodexFullRowpSI    & \CodexFullModalSI   \\
Codex seed 2            & \CodexFullTrainSII  & \CodexFullRowpSII   & \CodexFullModalSII  \\
Codex seed 3            & \CodexFullTrainSIII & \CodexFullRowpSIII  & \CodexFullModalSIII \\
\midrule
Gemini (best of 5)      & 5557.1 & 4272.7  & SLSQP \\
Claude run 1            & 5559.5 & 4264.9  & SGD-wrapper \\
500-start baseline      & 5540.7 & 4246.7  & ---  \\
\bottomrule
\end{tabular}
\end{sc}
\end{small}
\end{center}
\vskip -0.1in
\end{table}

[PLACEHOLDER: 1--2 paragraphs describing the Codex full-optimizer
results. Key questions: did Codex rediscover SLSQP (matching Gemini),
wrap topfarm\_sgd\_solve (matching Claude), or invent a third
approach? Does the broad-interface modal choice flip across seeds?]

\subsection{Narrow vs broad summary}

\begin{table}[t]
\caption{Per-agent narrow-vs-broad outcome, including Codex.
Discovery = novel schedule structure;
Reimplementation = community-standard solver wrapper.}
\label{tab:nbsummary}
\vskip 0.1in
\begin{center}
\begin{small}
\begin{sc}
\begin{tabular}{lll}
\toprule
Agent   & Narrow            & Broad                  \\
\midrule
Claude  & discovery (bumps) & reimpl (SGD wrapper)   \\
Gemini  & discovery (cos)   & reimpl (SLSQP)         \\
Codex   & \CodexNarrowVerdict & \CodexBroadVerdict   \\
\bottomrule
\end{tabular}
\end{sc}
\end{small}
\end{center}
\vskip -0.1in
\end{table}

[PLACEHOLDER: short interpretation of \cref{tab:nbsummary}.]

\section{Discussion}

[PLACEHOLDER: 1 paragraph on whether Codex strengthens or weakens the
companion paper's central claim. Three possible outcomes:
\begin{itemize}
\item All three agents discover-narrow / reimplement-broad: the claim
strengthens substantially, becomes "frontier-agent-invariant".
\item Codex breaks the pattern (e.g., reimplements under narrow):
the claim narrows to "Claude-and-Gemini-specific" and we should
investigate why Codex behaves differently.
\item Mixed: Codex matches on one mode but not the other; report
which specific transition fails.
\end{itemize}]

\paragraph{Limitations.}
$n{=}3$ agent seeds per cell; Codex agent-run variance not measured
beyond the 3 included runs. Single training farm, single held-out
polygon (inherited from the companion paper).
The structural classification is regex-based; a human review of all
six scripts is provided in the supplement to verify motifs.

\section{Conclusion}

[PLACEHOLDER: 2--3 sentences on whether Codex extends the FunWake
narrow-interface discovery result to a third frontier provider, and
what the 2026 frontier-agent landscape implies for using narrow
interfaces as a general-purpose algorithm-discovery elicitation
mechanism.]

\section*{Impact Statement}

This paper presents work whose goal is to advance the field of
Machine Learning. There are many potential societal consequences of
our work, none which we feel must be specifically highlighted here.

\bibliographystyle{icml2026}
\bibliography{references}

\end{document}
